% vim: ai sw=2 spell spelllang=cs

\shorthandoff{-} % cline fix

\begin{frame}
  \frametitle{Ladění jádra}

  \heading{LDD3 kap. 4} (zastaralá)

  \bigskip

  \heading{Jak zjistit chybu v jádře}
  \begin{itemize}
    \item Ladicími výpisy
    \pause
    \item Výstupem přes soubor (obsah velkých bufferů apod.)
    \pause
    \item Pád $\Rightarrow$ analýza Oops
    \pause
    \item Regrese $\Rightarrow$ \cmd{git bisect}
    \pause
    \item Automatické nástroje
    \pause
    \item (Debugger -- \cmd{kgdb} a další)
    \pause
    \item \ldots
  \end{itemize}
\end{frame}

\iffimuni
\section{Ladění jádra}
\frame{\sectionpage}
\fi

\begin{frame}[fragile]{Výpisy}

  \begin{itemize}
    \item \code{printk} již známe

    \pause

    \item Před formátovací řetězec se vkládá \code{KERN_*}
    \begin{itemize}
      \item Řetězec ve formě \code{<cislo>} nebo \code{\\001cislo}

      \pause

      \item Sděluje úroveň (důležitost) zprávy
      \item Definované v \header{linux/kernel.h}
    \end{itemize}

    \pause

    \item Zkratky: \code{pr_err}, \code{pr_info}, \dots
  \end{itemize}

  \onslide<2->

  \begin{block}{Příklad}
    \begin{lstlisting}
printk(KERN_ERR "button.c: Not enough memory\n");
printk(KERN_INFO "SMP mode deactivated.\n");
    \end{lstlisting}

    \onslide<4->

    \begin{lstlisting}
pr_err("button.c: Not enough memory\n");
pr_info("SMP mode deactivated.\n");
    \end{lstlisting}
  \end{block}

\end{frame}

\begin{frame}{Nastavení výpisů}

  \heading{Úroveň ovlivňuje, co se kde zobrazí}
  \begin{itemize}
    \item \cmd{dmesg} nastavit nelze
    \begin{itemize}
      \item Obsahuje vždy všechno
    \end{itemize}

    \pause

    \item Nastavení \file{/var/log/*} (syslogu) závislé na systému
  \end{itemize}

  \onslide<1->

  \begin{block}{Obvyklé nastavení}
  \begin{center}
  \begin{tabular}{|l|c|c|c|} \hline
  \multicolumn{1}{|c|}{\multirow{2}{*}{Důležitost}} &
  	\multicolumn{3}{c|}{Cíl} \\ \cline{2-4}
                    & \cmd{dmesg} & \file{/var/log/*}   & Konzole \\ \hline\hline
  \code{KERN_EMERG} & \cmark      & \cmark              & \cmark  \\ \hline
  \code{KERN_INFO}  & \cmark      & \cmark              & \xmark  \\ \hline
  \code{KERN_DEBUG} & \cmark      & \xmark              & \xmark  \\ \hline\hline
  \onslide<3->{Zapisuje:}
	& \onslide<3->{jádro}
	& \onslide<3->{syslog*}
	& \onslide<3->{jádro} \\ \hline
  \end{tabular}
  \end{center}
  \end{block}

\end{frame}

\begin{frame}{Nastavení konzole}

  \begin{itemize}
    \item Příkazem: \cmd{dmesg -n cislo}

    \pause

    \item Zprávy s vyšší úrovní se objeví na konzoli
    \begin{itemize}
      \item Vyšší úroveň = nižší číslo
    \end{itemize}

    \pause

    \item Hodnota 8 (debug) $\Rightarrow$ všechno dění
    \begin{itemize}
      \item Např.\ při ,,hard-locku''
    \end{itemize}
  \end{itemize}

  \medskip
  \pause

  \heading{Demo:} \file{pb173/03} + \cmd{dmesg -n}

\end{frame}

\iffimuni

\begin{frame}{Úkol}

  \heading{Práce s úrovněmi}
    \begin{enumerate}
    \item Přidejte si do kódu modulu \code{WARN_ON(1)}
    \item Zjistěte s jakou úrovní se varování vypisuje
    \begin{itemize}
      \item Pomocí postupného nastavování \cmd{dmesg -n}
    \end{itemize}
  \end{enumerate}

\end{frame}

\fi

\begin{frame}[fragile]
  \frametitle{Úrověň výpisů}

  \heading{Volba úrovně}
  \begin{itemize}
    \item Nechceme 10G \file{/var/log/messages}
    \begin{itemize}
      \item Ladicí výpisy: \code{KERN_DEBUG}
    \end{itemize}

    \pause

    \item Nechceme spoustu hlášek na konzoli
    \begin{itemize}
      \item Obvyklé chyby: \code{KERN_WARNING}, \code{KERN_ERR}
    \end{itemize}
  \end{itemize}

  \pause

  \begin{block}{Příklady}
  \begin{lstlisting}
printk(KERN_EMERG "Thread overran stack, or stack corrupted\n");
printk(KERN_ALERT "Unhandled fault: %s (0x%03x) at 0x%08lx\n", ...
printk(KERN_CRIT "Invalid DMA channel for ata\n");
printk(KERN_ERR "button.c: Not enough memory\n");
printk(KERN_WARNING "Only using first memory bank\n");
printk(KERN_NOTICE "PCI: Starting initialization.\n");
printk(KERN_INFO "SMP mode deactivated.\n");
printk(KERN_DEBUG "cosa%d: microcode started\n", ...
  \end{lstlisting} %stopzone
  \end{block}

\end{frame}

\begin{frame}
  \frametitle{Cíl výpisů}

  \heading{Co vše může být konzole pro výpisy}
  \begin{itemize}
    \item \file{tty0-tty*} (co když stroj umře?)

    \pause

    \item Sériová konzole (\file{ttyS*})
    \item Síťová konzole (\texttt{netconsole})

    \pause

    \item Všechno, co se zaregistruje pomocí \code{register_console}
  \end{itemize}

  \medskip
  \pause

  Nastavení pomocí parametru jádra (modulu)
  \begin{itemize}
    \item \texttt{console=tty1}
    \item \texttt{console=ttyS0}
    \item \texttt{netconsole=...}
  \end{itemize}

  \medskip
  \pause

  \doc{Documentation/console/console.txt}\\
  \doc{Documentation/serial-console.txt}

\end{frame}

\begin{frame}{Úkol}

  \heading{Použití netconsole ve virtuálním stroji}
  \begin{enumerate}
    \item Na hostitelovi: \cmd{nc -ul 0.0.0.0 60000}
    \item \cmd{dmesg -n 8} (vypisovat vše)
    \item \cmd{modinfo netconsole}
    \item \cmd{modprobe netconsole netconsole=@/ens3,60000@10.0.2.2/}
    \item Např.\ \cmd{echo h >/proc/sysrq-trigger}
  \end{enumerate}

  \smallskip

  Více informací v \doc{Documentation/networking/netconsole.txt}

\end{frame}

\begin{frame}[fragile]
  \frametitle{Nevýhody výpisů}

  \begin{itemize}
    \item Mnoho dat
    \begin{itemize}
      \item Znepřehlednění (typické pro ACPI, iwlwifi)
      \item Nestíhá se posílat na konzoli (ztráta informací)

      \pause

      \item \emph{Částečné řešení} (\header{linux/printk.h}):
      \begin{lstlisting}
if (printk_ratelimit())
  printk(...);
      \end{lstlisting}
    \end{itemize}

    \pause

    \item Binární data
    \begin{itemize}
      \item \cmd{dmesg} zobrazuje jen textová data

      \pause

      \item \emph{Částečné řešení} (\header{linux/kernel.h}):
      \begin{lstlisting}
print_hex_dump_bytes(...) // Vyzaduje DEBUG
print_hex_dump(...)
      \end{lstlisting}
    \end{itemize}
  \end{itemize}

  \pause

  \begin{center}
    \heading{Mnoho binárních dat se předává pomocí speciálních souborů}
  \end{center}

\end{frame}

\mysection{debugfs}

\begin{frame}
  \frametitle{debugfs}

  \begin{itemize}
    \item Předávání ladicích dat přes soubor

    \item Většinou v \file{/sys/kernel/debug/}
    \begin{itemize}
      \item Pokud ne: \cmd{mount -t debugfs none /sys/kernel/debug/}
    \end{itemize}

    \item \doc{Documentation/filesystems/debugfs.txt}
    \item \header{linux/debugfs.h}

    \item Vytvoření
    \begin{itemize}
      \item Adresář: \code{debugfs_create_dir}
      \item Symbolický odkaz: \code{debugfs_create_symlink}
      \item Vystavit proměnnou: \code{debugfs_create_u\{8,16,32,bool\}}

      \pause

      \item Celé \code{file_operations}: \code{debugfs_create_file}
      \begin{itemize}
	\item Lze si předat \code{data}, jsou pak v \code{open} v
	  \code{inode->i_private}
      \end{itemize}
    \end{itemize}

    \pause

    \item Smazání
    \begin{itemize}
      \item 1 položka: \code{debugfs_remove}
      \item Celý podstrom: \code{debugfs_remove_recursive}
    \end{itemize}
  \end{itemize}

\end{frame}

\begin{frame}{Úkol}

  \heading{Vytvoření položek v debugfs}
  \begin{enumerate}
  \item Vytvořte adresář (\code{debugfs_create_dir})

  \item Vytvořte v něm soubor (\code{debugfs_create_u16})
  \begin{itemize}
    \item Vrací nějakou hodnotu (např.\ 11320)
  \end{itemize}

  \item Přeložte, vložte do systému (\cmd{make}, \cmd{insmod})

  \item Přečtěte soubor
  \begin{itemize}
    \item \cmd{cat /sys/kernel/debug/<dir>/<file>}
  \end{itemize}

  \end{enumerate}

\end{frame}

\mysection{sysrq + strace}

\begin{frame}
  \frametitle{SysRq}

  \begin{itemize}
    \item System-request (\doc{Documentation/sysrq.txt})
    \begin{itemize}
      \item Klávesnicí: Alt-PrintScreen
      \item Vzdáleně: \file{/proc/sysrq-trigger}
    \end{itemize}

    \pause

    \item Funguje (většinou) i po pádu systému
    \item Zapnutí: \cmd{sysctl kernel.sysrq=1}
  \end{itemize}

  \pause

  \begin{block}{Příklady SysRq}
  \begin{center}
  \small
  \begin{tabular}{|l|l|} \hline
  h & zkrácená nápověda \\ \hline
  r & ,,reset klávesnice'' \\ \hline
  p & výpis tras procesů na CPU \\ \hline
  t & výpis tras všech procesů \\ \hline
  w & výpis blokujících procesů \\ \hline
  s-u-s-b & sync + unmount + reboot \\ \hline
  \end{tabular} \\[2mm]
  \end{center}
  \end{block}

\end{frame}


\begin{frame}{Úkol}

  \heading{Sysrq}
  \begin{enumerate}
    \item Vypište si nápovědu (\cmd{sysrq-h})
    \item Vypište si trasy procesů na CPU (\cmd{sysrq-p})
    \item Vypište si informace o paměti (\cmd{sysrq-m})
  \end{enumerate}

  \bigskip

  Pozn.: v Qemu lze vyzkoušet v ovládací konzoli (Ctrl-Alt-2) příkazem
  \texttt{sendkey alt-sysrq-písmeno}

\end{frame}

\iffimuni
\begin{frame}
  \frametitle{strace}

  \begin{itemize}
  \item Chci-li znát komunikaci jádro $\leftrightarrow$ program
  \item \doc{man strace}
  \item Zachytává syscally
  \item Dekóduje a vypisuje jejich parametry na STDOUT
  \end{itemize}

  \heading{Demo:} \cmd{strace -e open,read cat .../debug/...}

\end{frame}

\begin{frame}[fragile]{Analýza Oops}

  \footnotesize
  \begin{lstlisting}
my_init(void) { char *ptr = (void *)5; *ptr = 5; }
  \end{lstlisting}
  \begin{lstlisting}{language=XML}
BUG: unable to handle kernel NULL pointer dereference at 00000000000005
IP: [<ffffffffa0039042>] my_init+0x12/0x30 [pb173]
PGD 3daff067 PUD 3c713067 PMD 0 
...
Modules linked in: pb173(+) mperf fuse ...
Pid: 2506, comm: insmod Not tainted 2.6.36-rc4-16-default #1 /Bochs
RIP: 0010:[<ffffffffa0039042>]  [<ffffffffa0039042>] my_init+0x12/0x30 [pb173]
RSP: 0018:ffff88003c741f28  EFLAGS: 00010292
RAX: 0000000000000017 RBX: ffffffffa0039180 RCX: 0000000000000e91
...
CR2: 0000000000000005 CR3: 000000003bb47000 CR4: 00000000000006f0
Process insmod (pid: 2506, ...)
Stack: ...
Call Trace:
 [<ffffffff810002da>] do_one_initcall+0x3a/0x170
 [<ffffffff8108f69a>] sys_init_module+0xba/0x210
 [<ffffffff81002efb>] system_call_fastpath+0x16/0x1b
 [<00007fc3386320ba>] 0x7fc3386320ba
Code: ... e8 63 c8 41 e1 <c6> 04 25 05 00 00 ...
  \end{lstlisting}

\end{frame}
\fi
